\documentclass[11pt]{article}
\usepackage{amsmath,amssymb}
\usepackage{geometry}
\geometry{margin=1in}

\begin{document}

\begin{center}
\textbf{CSE400: Fundamentals of Probability in Computing}\\
\textbf{Lecture 6: Discrete Random Variables, Expectation and Problem Solving}\\
Dhaval Patel, PhD \\
January 22, 2025
\end{center}

\section*{1. Random Variables}

\subsection*{Definition}
Let $\Omega$ be a sample space.  
A \emph{random variable} $X$ on $\Omega$ is a function
\[
X : \Omega \rightarrow \mathbb{R}
\]
that assigns a real number $X(\omega)$ to each sample point $\omega \in \Omega$. :contentReference[oaicite:0]{index=0}

\subsection*{Discrete Random Variables}
In this lecture, we restrict attention to \emph{discrete} random variables.  
A random variable is said to be discrete if it can take values in a set that is either finite or countably infinite.  
Although $X$ maps into $\mathbb{R}$, the set of values $\{X(\omega) : \omega \in \Omega\}$ is a discrete subset of $\mathbb{R}$. :contentReference[oaicite:1]{index=1}

\subsection*{Visualization}
The distribution of a discrete random variable can be visualized using a bar diagram:
\begin{itemize}
\item The $x$-axis represents the possible values of the random variable.
\item The height of the bar at value $a$ is $\Pr[X = a]$.
\end{itemize}
Each probability is computed from the probability of the corresponding event in the sample space. :contentReference[oaicite:2]{index=2}

\section*{2. Discrete vs Continuous Random Variables}

\subsection*{Discrete Random Variable}
A discrete random variable has:
\begin{itemize}
\item Countable support,
\item A probability mass function (PMF),
\item Probabilities assigned to single values,
\item Strictly positive probabilities for each possible value.
\end{itemize} :contentReference[oaicite:3]{index=3}

\subsection*{Continuous Random Variable}
A continuous random variable has:
\begin{itemize}
\item Uncountable support,
\item A probability density function (PDF),
\item Probabilities assigned to intervals of values,
\item Zero probability for each individual value.
\end{itemize} :contentReference[oaicite:4]{index=4}

\section*{3. Example: Tossing Three Fair Coins}

\subsection*{Experiment Description}
Suppose an experiment consists of tossing three fair coins.  
Let $Y$ denote the number of heads observed.

Possible values:
\[
Y \in \{0,1,2,3\}.
\]

\subsection*{Probability Distribution}
\[
\Pr(Y=0)=\Pr(t,t,t)=\frac{1}{8},
\]
\[
\Pr(Y=1)=\Pr(t,t,h \text{ or } t,h,t \text{ or } h,t,t)=\frac{3}{8},
\]
\[
\Pr(Y=2)=\Pr(t,h,h \text{ or } h,t,h \text{ or } h,h,t)=\frac{3}{8},
\]
\[
\Pr(Y=3)=\Pr(h,h,h)=\frac{1}{8}.
\]

Since $Y$ must take one of the values $0,1,2,3$, we verify normalization:
\[
\sum_{i=0}^{3} \Pr(Y=i) = 1.
\] 
:contentReference[oaicite:5]{index=5}

\section*{4. Probability Mass Function (PMF)}

\subsection*{Definition}
A random variable that can take at most a countable number of possible values is said to be discrete.

Let $X$ be a discrete random variable with range
\[
R_X = \{x_1, x_2, x_3, \dots\}
\]
(finite or countably infinite).  
The function
\[
P_X(x_k) = \Pr(X = x_k), \quad k = 1,2,3,\dots
\]
is called the \emph{Probability Mass Function (PMF)} of $X$. :contentReference[oaicite:6]{index=6}

\subsection*{Property}
Since $X$ must take one of the values $x_k$,
\[
\sum_{k=1}^{\infty} P_X(x_k) = 1.
\] 
:contentReference[oaicite:7]{index=7}

\section*{5. PMF Example}

\subsection*{Given}
The PMF of a random variable $X$ is defined by
\[
p(i) = c \frac{\lambda^i}{i!}, \quad i=0,1,2,\dots
\]
where $\lambda > 0$.

\subsection*{Finding the Constant $c$}
Using normalization:
\[
\sum_{i=0}^{\infty} p(i) = 1
\]
\[
c \sum_{i=0}^{\infty} \frac{\lambda^i}{i!} = 1.
\]
Since
\[
\sum_{i=0}^{\infty} \frac{\lambda^i}{i!} = e^{\lambda},
\]
we obtain
\[
c e^{\lambda} = 1 \quad \Rightarrow \quad c = e^{-\lambda}.
\]

\subsection*{Computations}
\[
\Pr(X=0) = e^{-\lambda}.
\]

\[
\Pr(X > 2) = 1 - \Pr(X \le 2)
= 1 - \Pr(X=0) - \Pr(X=1) - \Pr(X=2).
\] 
:contentReference[oaicite:8]{index=8}

\section*{6. Bayes' Theorem}

\subsection*{Bayes Formula}
Using
\[
\Pr(A \cap B_i) = \Pr(B_i \mid A)\Pr(A),
\]
we obtain
\[
\Pr(B_i \mid A) = \frac{\Pr(A \mid B_i)\Pr(B_i)}{\sum_{j=1}^{n} \Pr(A \mid B_j)\Pr(B_j)}.
\]

This is known as the \emph{Bayes Formula (Proposition 3.1)}. :contentReference[oaicite:9]{index=9}

\subsection*{Terminology}
\begin{itemize}
\item $\Pr(B_i)$: \emph{a priori probability}.
\item $\Pr(B_i \mid A)$: \emph{posterior probability}.
\end{itemize}

\section*{7. Bayes' Theorem Example: Auditorium with 30 Rows}

\subsection*{Problem Description}
An auditorium has 30 rows of seats.
\begin{itemize}
\item Row 1 has 11 seats,
\item Row 2 has 12 seats,
\item Row 3 has 13 seats,
\item $\vdots$
\item Row 30 has 40 seats.
\end{itemize}

A door prize is given by randomly selecting:
\begin{enumerate}
\item a row (each row equally likely),
\item then a seat within that row (each seat equally likely).
\end{enumerate}

\subsection*{Computations}
Let $S_{15}$ denote the event that seat 15 is selected.  
Let $R_{20}$ denote the event that row 20 is selected.

\[
\Pr(S_{15} \mid R_{20}) = \frac{1}{30}.
\]

Total probability:
\[
\Pr(S_{15}) = \sum_{k=1}^{30} \Pr(S_{15} \mid R_k)\Pr(R_k).
\]

Using Bayes' formula:
\[
\Pr(R_{20} \mid S_{15})
= \frac{\Pr(S_{15} \mid R_{20})\Pr(R_{20})}{\Pr(S_{15})}.
\] 
:contentReference[oaicite:10]{index=10}

\end{document}
